%%%%%%%%%%%%%%%%%%%%%%%%%%%%%%%%%%%%%%%%%%%%%%%%%%%%%%%%%%%%%%%%%%%%%%%%%%%%%%%
%% Seção 2 --- Fundamentação teórica e trabalhos relacionados
%%
%% Junta, por decisão do usuário, a fundamentação e a revisão da literatura, antes
%% separadas: cada tema traz o conceito e o que outros autores publicaram sobre ele.
%% O projeto de hardware (invólucro, térmica, bateria) é trabalho futuro e não entra.
%%
%% Os valores do dispositivo vêm de:
%%   rolamento e rotação   docs/relatorios/relatorio-kaelix.tex, tab:parametros
%%   enlace LoRa           src/comms/lora_config.h (SF9, 125 kHz, 4/5, preâmbulo 8)
%%   carga útil            src/comms/lora.h (21 bytes)
%%   MPU6050               folha de dados rev. 3.1 e mapa de registradores
%%                         rev. 4.0, conferidos no documento primário
%%%%%%%%%%%%%%%%%%%%%%%%%%%%%%%%%%%%%%%%%%%%%%%%%%%%%%%%%%%%%%%%%%%%%%%%%%%%%%%

\makeatletter
\@ifundefined{toprule}{%
  \newcommand{\toprule}{\hline}%
  \newcommand{\midrule}{\hline}%
  \newcommand{\bottomrule}{\hline}%
}{}
\@ifundefined{emandamento}{%
  \@ifundefined{textcolor}%
    {\newcommand{\emandamento}[1]{#1}}%
    {\newcommand{\emandamento}[1]{\textcolor{red}{#1}}}%
}{}
\makeatother

\chapter{Fundamentação teórica e trabalhos relacionados}
\label{cap:fundamentacao}

Esta seção apresenta os conceitos usados no restante do texto e o que outros autores publicaram
sobre cada tema, na ordem em que a informação percorre o dispositivo: a vibração da máquina e as
assinaturas das falhas, a aquisição do sinal, a avaliação pela norma, o detector de anomalias e o
nó sensor com seu enlace de rádio. Sempre que possível, os conceitos são aplicados ao Kaelix, para
que os números das seções seguintes possam ser refeitos. A seção termina com as lacunas encontradas
na literatura.

A revisão partiu das escolhas de projeto e dos resultados obtidos e procurou, para cada um, trabalhos
que os confirmassem, os contradissessem ou não tratassem deles. \emandamento{Informar aqui as bases
consultadas, os termos de busca e o período da busca.} Os autores, os títulos e os identificadores
foram conferidos no registro do editor. Quando o texto integral não pôde ser acessado, isso consta em
nota na própria referência. \emandamento{Antes da versão final, conferir o texto das obras citadas e
as páginas de Prodanov e Freitas.}

%%%%%%%%%%%%%%%%%%%%%%%%%%%%%%%%%%%%%%%%%%%%%%%%%%%%%%%%%%%%%%%%%%%%%%%%%%%%%%%
\section{Manutenção de máquinas rotativas}
\label{sec:fund-manutencao}

A manutenção de equipamentos industriais é usualmente classificada em três estratégias, conforme o
momento da intervenção \cite{mobley2002}. Na manutenção \emph{corretiva}, o equipamento opera até
falhar e só então é reparado. Não há custo de planejamento, mas a falha ocorre em momento não
escolhido e pode danificar componentes íntegros. Na manutenção \emph{preventiva}, a intervenção
ocorre em intervalos fixos de tempo ou de uso, definidos pela vida média esperada dos componentes. A
falha inesperada fica menos frequente, mas componentes em bom estado são trocados sem necessidade, e
falhas anteriores ao intervalo continuam sem detecção. Na manutenção \emph{preditiva}, ou baseada em
condição, a intervenção é decidida pelo estado medido do equipamento: acompanha-se uma grandeza
física que reflete sua condição e intervém-se quando ela indica degradação.

\citeonline{jardine2006} descrevem um programa de manutenção baseada em condição em três etapas:
aquisição de dados, processamento dos dados e tomada de decisão. Na aquisição, escolhem-se a grandeza,
o sensor, o ponto de medição e a taxa. No processamento, extraem-se indicadores dos dados brutos. Na
decisão, esses indicadores definem se a máquina precisa de intervenção.

Entre as grandezas usadas no monitoramento de máquinas rotativas, a vibração é a mais difundida,
porque a maior parte dos defeitos mecânicos (desbalanceamento, desalinhamento, folga, desgaste de
rolamentos e de engrenagens) altera as forças dinâmicas no interior da máquina, e essas forças se
transmitem à carcaça, onde podem ser medidas sem interromper a operação \cite{randall2011livro}.

%%%%%%%%%%%%%%%%%%%%%%%%%%%%%%%%%%%%%%%%%%%%%%%%%%%%%%%%%%%%%%%%%%%%%%%%%%%%%%%
\section{Vibração e assinaturas de falha}
\label{sec:fund-vibracao}

A vibração de uma máquina é descrita por grandezas ligadas por derivação e resumida em descritores
escalares. Cada tipo de defeito altera essa vibração de um modo característico, e essas assinaturas
dividem-se em dois grupos, conforme a faixa de frequência em que aparecem.

\subsection{Deslocamento, velocidade e aceleração}

Vibração é o movimento oscilatório de um ponto em torno de uma posição de equilíbrio. O movimento pode
ser descrito por três grandezas relacionadas: o deslocamento $x(t)$, a velocidade
$v(t) = \mathrm{d}x/\mathrm{d}t$ e a aceleração $a(t) = \mathrm{d}v/\mathrm{d}t$. Para um movimento
senoidal de frequência $f$ e amplitude de deslocamento $X$, com $\omega = 2\pi f$,

\begin{equation}
x(t) = X \sin(\omega t), \qquad
v(t) = \omega X \cos(\omega t), \qquad
a(t) = -\omega^2 X \sin(\omega t).
\label{eq:fund-senoide}
\end{equation}

As amplitudes das três grandezas diferem, portanto, por fatores de $\omega$. Para um mesmo deslocamento,
a velocidade cresce proporcionalmente à frequência, e a aceleração, ao quadrado dela. Por isso, a grandeza
escolhida muda o peso de cada frequência: o deslocamento realça as baixas frequências, a aceleração realça
as altas, e a velocidade fica entre as duas. Conhecida uma das grandezas, as outras podem ser obtidas por
derivação ou integração. Um acelerômetro, que mede $a(t)$, pode fornecer a velocidade por integração,
operação discutida na seção~\ref{sec:fund-integracao}.

\subsection{Descritores escalares de um sinal}
\label{sec:fund-descritores}

Um sinal de vibração registrado durante um intervalo é frequentemente resumido por descritores escalares.
Para uma janela de $N$ amostras $x_n$, com média $\bar{x}$ e desvio-padrão $\sigma$, os mais usados são o
valor eficaz, o pico, o pico a pico, o fator de crista (FC) e a curtose \cite{randall2011livro}:

\begin{equation}
x_{\mathrm{rms}} = \sqrt{\frac{1}{N}\sum_{n=1}^{N} x_n^2},
\qquad
x_{\mathrm{pico}} = \max_n |x_n|,
\qquad
\mathrm{FC} = \frac{x_{\mathrm{pico}}}{x_{\mathrm{rms}}},
\qquad
\kappa = \frac{\frac{1}{N}\sum_{n=1}^{N} (x_n-\bar{x})^4}{\sigma^4}.
\label{eq:fund-descritores}
\end{equation}

O valor eficaz ($x_{\mathrm{rms}}$) mede a energia média do sinal e é a base da avaliação de severidade
(seção~\ref{sec:fund-norma}). O valor de pico e o pico a pico (diferença entre o maior e o menor valor da
janela) medem a excursão máxima. O fator de crista e a curtose ($\kappa$) medem a \emph{impulsividade}:
quanto o sinal concentra sua energia em eventos breves e intensos. Uma senoide pura tem
$\mathrm{FC} = \sqrt{2} \approx \num{1,41}$ e $\kappa = \num{1,5}$, e um ruído gaussiano tem $\kappa = 3$.
O \emph{firmware} usa a curtose em excesso, $\kappa - 3$, que vale 0 para esse ruído. Um sinal composto por
impactos espaçados, como o de um defeito localizado em rolamento (seção~\ref{sec:fund-rolamento}), tem fator
de crista e curtose bem maiores, porque poucos valores extremos dominam a média de quarta potência. Esses
descritores não indicam em que frequência a energia está nem a relação entre componentes de frequências
diferentes. Essa informação vem da análise espectral (seção~\ref{sec:fund-espectro}).

\subsection{Defeitos síncronos com a rotação}

Desbalanceamento e desalinhamento produzem forças que se repetem a cada volta do eixo e, por isso, aparecem
em múltiplos inteiros da frequência de rotação $f_r$, as \emph{harmônicas} da rotação
\cite{randall2011livro}. No desbalanceamento, o centro de massa do rotor não coincide com o eixo de rotação,
e a força centrífuga resultante gira com o eixo. A assinatura dominante é a componente em $1\times f_r$, cuja
amplitude cresce com a massa desbalanceada. No desalinhamento, os eixos acoplados não são colineares, e o
acoplamento impõe esforços que variam mais de uma vez por volta. A assinatura típica é o aumento da componente
em $2\times f_r$ em relação à de $1\times f_r$.

Como as rotações industriais usuais estão entre algumas centenas e alguns milhares de rotações por minuto,
essas componentes ficam entre poucos hertz e algumas centenas de hertz. Um motor a \SI{1750}{r/min}, por
exemplo, tem $f_r = \SI{29,17}{\hertz}$, e suas três primeiras harmônicas estão abaixo de \SI{100}{\hertz}.

\subsection{Defeitos em rolamentos}
\label{sec:fund-rolamento}

O rolamento de esferas é composto por uma pista interna, solidária ao eixo, uma pista externa, fixa à carcaça,
e um conjunto de esferas mantidas equidistantes por uma gaiola. Um defeito localizado, como uma trinca ou um
lascamento em uma das pistas, é atingido periodicamente pelas esferas, e cada passagem produz um impacto. A taxa
de repetição desses impactos depende apenas da geometria do rolamento e da rotação. Ela é chamada frequência de
passagem de esferas, na pista externa (BPFO) ou na pista interna (BPFI), conforme a posição do defeito, e vale
\cite{randall2011}

\begin{equation}
\mathrm{BPFO} = \frac{N_b}{2}\, f_r \left(1 - \frac{d}{D_p}\cos\theta\right),
\qquad
\mathrm{BPFI} = \frac{N_b}{2}\, f_r \left(1 + \frac{d}{D_p}\cos\theta\right),
\label{eq:fund-bpfo}
\end{equation}

\noindent
com $N_b$ o número de esferas, $d$ o diâmetro da esfera, $D_p$ o diâmetro primitivo e $\theta$ o ângulo de
contato. Para o rolamento SKF~6205 ($N_b = 9$, $d/D_p = \num{0,203}$, $\theta = 0$) a \SI{1750}{r/min}, as
expressões dão $\mathrm{BPFO} = \SI{104,6}{\hertz}$ e $\mathrm{BPFI} = \SI{157,9}{\hertz}$.

A maior parte da energia do defeito, porém, fica fora da frequência de repetição. Cada impacto é um evento breve,
que excita as ressonâncias da estrutura entre o rolamento e o ponto de medição, tipicamente na faixa de alguns
quilohertz. O sinal resultante é uma sequência de oscilações amortecidas na frequência de ressonância, que se
repetem na taxa BPFO ou BPFI. No espectro, a energia do defeito fica em torno da ressonância, e a frequência
característica aparece como \emph{modulação} dessa portadora de alta frequência \cite{randall2011}. Nos estágios
iniciais do defeito, as linhas na própria BPFO tendem a ser fracas diante das demais componentes da máquina.

Por isso, o procedimento usual de diagnóstico de rolamentos é a \emph{análise de envelope}. O sinal é filtrado em
uma banda em torno da ressonância, calcula-se seu envelope, isto é, a amplitude instantânea, pelo módulo do sinal
analítico obtido com a transformada de Hilbert, e o espectro desse envelope mostra linhas na frequência de repetição
dos impactos \cite{randall2011}. A curtose espectral, isto é, a curtose calculada em função da frequência, é a
ferramenta usada para localizar a banda em que o sinal é mais impulsivo \cite{antoni2007,antoni2006}. Esse
procedimento só funciona se a banda de ressonância tiver sido adquirida.

%%%%%%%%%%%%%%%%%%%%%%%%%%%%%%%%%%%%%%%%%%%%%%%%%%%%%%%%%%%%%%%%%%%%%%%%%%%%%%%
\section{Aquisição do sinal}
\label{sec:fund-sinais}

A taxa de amostragem fixa a maior frequência representável, a duração da janela fixa a resolução espectral e o
transdutor determina a banda e o piso de ruído da medição.

\subsection{Amostragem, frequência de Nyquist e rebatimento}
\label{sec:fund-amostragem}

Um sistema digital registra amostras $a_n = a(n/f_s)$ do sinal contínuo $a(t)$, tomadas à taxa de amostragem
$f_s$. O teorema da amostragem estabelece que um sinal pode ser reconstruído com base em suas amostras se não
contiver componentes acima de $f_s/2$, a \emph{frequência de Nyquist} \cite{brandt2011}. Uma componente acima
desse limite não desaparece na amostragem: suas amostras são idênticas às de uma componente de frequência mais
baixa, dentro da faixa $[0, f_s/2]$. O fenômeno chama-se \emph{rebatimento} (\emph{aliasing}), e a frequência
aparente de uma componente em $f$ é

\begin{equation}
f_{\mathrm{ap}} = \left| f - f_s \cdot \mathrm{round}\!\left(\frac{f}{f_s}\right) \right|.
\label{eq:fund-alias}
\end{equation}

Depois da amostragem, a componente rebatida não pode mais ser separada de uma componente real naquela frequência.
Para evitar isso, é preciso remover o conteúdo acima de $f_s/2$ \emph{antes} da amostragem, com um filtro
passa-baixas chamado filtro \emph{antirrebatimento} (\emph{anti-aliasing}). Pela mesma razão, reduzir a taxa de um
sinal já digitalizado, operação chamada \emph{decimação}, exige filtrar o sinal abaixo da nova frequência de Nyquist
antes de descartar amostras \cite{brandt2011}.

No Kaelix, com $f_s = \SI{1}{\kilo\hertz}$, a frequência de Nyquist é \SI{500}{\hertz}. Uma ressonância de rolamento
em \SI{3,5}{\kilo\hertz} que chegasse sem filtro ao conversor seria observada, pela Equação~\eqref{eq:fund-alias}, em
$|3500 - 1000 \cdot 4| = \SI{500}{\hertz}$, onde não há nenhuma componente física.

\subsection{Análise espectral}
\label{sec:fund-espectro}

A representação de um sinal em função da frequência é obtida pela transformada discreta de Fourier (DFT), calculada
na prática pelo algoritmo da transformada rápida (FFT). Para uma janela de $N$ amostras,

\begin{equation}
X_k = \sum_{n=0}^{N-1} x_n\, e^{-\mathrm{j}2\pi k n / N}, \qquad k = 0, 1, \dots, N-1,
\label{eq:fund-dft}
\end{equation}

\noindent
e o coeficiente $X_k$ corresponde à frequência $k\,f_s/N$ \cite{brandt2011}. A \emph{resolução} em frequência é
$\Delta f = f_s/N$, o inverso da duração da janela, e duas componentes mais próximas que $\Delta f$ não são separadas.
A janela finita também provoca \emph{vazamento espectral}: uma componente cuja frequência não coincide com um múltiplo
de $\Delta f$ espalha energia pelos coeficientes vizinhos. Esse efeito diminui quando o sinal é multiplicado por uma
função de janela que o leve suavemente a zero nas bordas, como a de Hann \cite{brandt2011}. Para sinais aleatórios, o
espectro de uma única janela é uma estimativa muito ruidosa, e usa-se a \emph{densidade espectral de potência}, obtida
pela média dos espectros de várias janelas.

No Kaelix, com $f_s = \SI{1}{\kilo\hertz}$ e janelas de \num{512} amostras, a janela dura \SI{0,512}{\second} e a
resolução é de \SI{1,95}{\hertz}, suficiente para separar as harmônicas de uma rotação de \SI{29}{\hertz}.

\subsection{Acelerômetros piezoelétricos e MEMS}

O transdutor de referência no monitoramento de máquinas é o acelerômetro piezoelétrico: uma massa sísmica apoiada sobre
um cristal que gera carga elétrica proporcional à força e, portanto, à aceleração. Sua frequência de ressonância interna
é elevada, o que lhe dá resposta plana até vários quilohertz, e seu ruído é baixo. Em contrapartida, ele exige
condicionamento de sinal dedicado e custa caro \cite{albarbar2008}.

Os acelerômetros MEMS capacitivos são fabricados pelos mesmos processos dos circuitos integrados. Uma microestrutura
suspensa por molas se desloca sob aceleração, e o deslocamento altera a capacitância entre eletrodos. O conversor
analógico-digital fica no mesmo encapsulamento, e o sensor entrega amostras digitais por barramento serial. O custo é de
uma a duas ordens de grandeza menor que o de um piezoelétrico, mas a banda e o ruído dependem muito do modelo. Há MEMS
projetados para monitoramento de condição, com banda de dezenas de quilohertz, e MEMS de consumo, projetados para orientação
e movimento humano, com banda de algumas centenas de hertz \cite{albarbar2008,fidali2024}. Para decidir se um MEMS serve
para uma aplicação de vibração, consideram-se a banda útil, a densidade espectral de ruído, a faixa de medição e a
estabilidade do zero.

A densidade espectral de ruído, expressa em $\si{\micro g}/\sqrt{\si{\hertz}}$, indica o ruído por unidade de banda. Para
ruído branco, o ruído eficaz na banda $B$ é a densidade multiplicada por $\sqrt{B}$. Quanto mais larga a banda analisada,
maior o ruído total e maior a menor amplitude de vibração que o sensor consegue distinguir.

\subsection{O acelerômetro MPU6050}
\label{sec:fund-mpu6050}

O Kaelix foi planejado para o MPU6050, uma unidade de medição inercial de consumo que combina acelerômetro e giroscópio de
três eixos. As características relevantes para medição de vibração, conforme a folha de dados e o mapa de registradores do
fabricante, estão na Tabela~\ref{tab:fund-mpu6050}.

\begin{table}[!htb]
\IBGEtab{%
\caption{Características do acelerômetro do MPU6050 relevantes para medição de vibração}%
\label{tab:fund-mpu6050}%
}{%
\begin{tabular}{@{}ll@{}}
\toprule
Característica & Valor \\
\midrule
Faixa de medição                         & $\pm2$, $\pm4$, $\pm8$ ou $\pm16$~g, programável \\
Resolução do conversor                   & \num{16} bits \\
Taxa de saída de dados                   & \SIrange{4}{1000}{\hertz}, programável \\
Filtro passa-baixas digital (DLPF)       & \SIrange{5}{260}{\hertz}, programável \\
Densidade espectral de ruído             & \SI{400}{\micro g}/$\sqrt{\si{\hertz}}$ a \SI{10}{\hertz} \\
Tolerância de zero (eixos X e Y)         & $\pm\SI{50}{\milli g}$ \\
\bottomrule
\end{tabular}%
}{%
\fonte{InvenSense (\citeyear{invensense2011ps}, \citeyear{invensense2012rm}).}%
}
\end{table}

Dois pontos da tabela afetam o uso do sensor em vibração. O primeiro é a banda. A taxa máxima de saída do acelerômetro é de
\SI{1}{\kilo\hertz}, o que põe a frequência de Nyquist em \SI{500}{\hertz}, mas quem fixa a banda efetiva é o filtro
passa-baixas digital interno, configurado pelo campo \texttt{DLPF\_CFG} (Tabela~\ref{tab:fund-dlpf}). No MPU6050, na configuração de banda
mais larga, \texttt{DLPF\_CFG}~=~0, que é também o valor após a inicialização, a banda do acelerômetro é de \SI{260}{\hertz},
e não há configuração em que o acelerômetro opere sem esse filtro \cite{invensense2012rm}. O \emph{firmware} usa
\texttt{DLPF\_CFG}~=~1, com banda de \SI{184}{\hertz}. O mapa de registradores informa a banda e o atraso de cada configuração,
mas não a ordem do filtro nem sua atenuação acima da banda, de modo que o fabricante não especifica a proteção contra
rebatimento.

\begin{table}[!htb]
\IBGEtab{%
\caption{Banda e atraso do filtro passa-baixas digital do acelerômetro do MPU6050, por configuração do campo \texttt{DLPF\_CFG}, com taxa de saída de \SI{1}{\kilo\hertz}}%
\label{tab:fund-dlpf}%
}{%
\begin{tabular}{@{}ccc@{}}
\toprule
\texttt{DLPF\_CFG} & Banda (\si{\hertz}) & Atraso (\si{\milli\second}) \\
\midrule
0 & 260 & \num{0,0}  \\
1 & 184 & \num{2,0}  \\
2 & 94  & \num{3,0}  \\
3 & 44  & \num{4,9}  \\
4 & 21  & \num{8,5}  \\
5 & 10  & \num{13,8} \\
6 & 5   & \num{19,0} \\
\bottomrule
\end{tabular}%
}{%
\fonte{\citeonline{invensense2012rm}.}%
}
\end{table}

O segundo ponto é o ruído. Tratando a densidade de \SI{400}{\micro g}/$\sqrt{\si{\hertz}}$ como ruído branco, o ruído eficaz na
banda de \SI{260}{\hertz} é de $400 \cdot \sqrt{260} \approx \SI{6,5}{\milli g}$, aproximadamente
\SI{0,063}{\meter\per\second\squared}. Na banda de \SI{184}{\hertz}, o mesmo cálculo dá \SI{5,4}{\milli g}. A tolerância de zero
de $\pm\SI{50}{\milli g}$ significa que o sensor pode apresentar uma componente contínua de até aproximadamente
\SI{0,5}{\meter\per\second\squared} sem estar fora de especificação. A integração no domínio do tempo acumula esse tipo de desvio
(seção~\ref{sec:fund-integracao}).

\subsection{Limites de banda e de ruído na literatura}
\label{sec:rev-aquisicao}

Os estudos de adequação de sensor apontam o limite superior de banda como a variável mais importante, e nenhum testou um dispositivo
próximo da faixa do Kaelix. A comparação de três acelerômetros MEMS capacitivos com um piezoelétrico de referência reprovou o de banda
mais estreita, de \SI{1500}{\hertz}, para monitoramento de condição \cite{albarbar2008}. Os grupos que diagnosticam rolamento com MEMS
escolhem sensores de dezenas de quilohertz \cite{landi2023,sharma2024}, e uma implementação de baixo custo usa dois acelerômetros: um
de banda estreita para desbalanceamento e outro de \SI{40}{\kilo\hertz} para rolamento \cite{jakobsen2024}. Os produtos comerciais
medem, ao menos em um eixo, uma ordem de grandeza acima do limite de uma IMU de consumo
\cite{bentlynevada2025,erbessd2025,tractian2022}.

O piso de ruído também limita a medição, independentemente da banda. A densidade de ruído da folha de dados de uma IMU de consumo,
integrada sobre a banda normativa, fica acima do ruído admissível para separar as zonas mais baixas da norma. Os estudos sobre MEMS
apontam a falta de calibração rastreável desses dispositivos \cite{albarbar2008,rodrigues2021}, e a avaliação de sensibilidade
diagnóstica de sensores digitais de catálogo chega à mesma conclusão \cite{fidali2024}. \emandamento{Não foi possível confirmar em
fonte primária a estimativa de ruído admissível usualmente atribuída à literatura de fabricante \cite{looney2017}, e a comparação com
a densidade medida ainda não foi feita.}

O filtro antirrebatimento é pouco discutido nessa literatura, embora afete dispositivos de consumo. A busca não encontrou, na literatura
revisada por pares sobre diagnóstico de rolamento, descrição experimental do efeito do rebatimento sobre descritores como curtose e fator
de crista. O efeito aparece apenas em notas de aplicação de fabricante, que não foram conferidas nesta revisão.

%%%%%%%%%%%%%%%%%%%%%%%%%%%%%%%%%%%%%%%%%%%%%%%%%%%%%%%%%%%%%%%%%%%%%%%%%%%%%%%
\section{Avaliação da severidade pela norma}
\label{sec:fund-norma}

A avaliação da vibração de máquinas industriais medida em partes não rotativas é normalizada pela ISO 20816-3:2022, que substituiu a
ISO 10816-3:2009 \citemesmoautor{International Organization for Standardization}{iso108163_2009}{iso208163_2022}. A norma aplica-se a
máquinas com potência acima de \SI{15}{\kilo\watt} e rotação entre \SI{120}{r/min} e \SI{30000}{r/min} e adota como grandeza de avaliação
a \emph{velocidade eficaz de banda larga}. Segundo o texto da norma, vibrações com o mesmo valor eficaz de velocidade na faixa de
\SIrange{10}{1000}{\hertz} podem, em geral, ser consideradas de mesma severidade, ao passo que critérios em deslocamento ou em aceleração
variariam com a frequência \cite{iso208163_2022}.

Pela mesma razão, a norma exige do instrumento resposta plana em \emph{ao menos} \SIrange{10}{1000}{\hertz}, com limite inferior de até
\SI{2}{\hertz} para máquinas com rotação próxima ou abaixo de \SI{600}{r/min}, e determina que o sinal de acelerômetros montados na carcaça
seja integrado para fornecer velocidade. Um instrumento com banda menor que a exigida não atende à norma, mesmo que o processamento posterior
esteja correto. A edição anterior delega a definição de resposta plana a uma norma de instrumentação, que usa tolerância de banda passante de
aproximadamente \SI{10}{\percent} entre \SI{12,6}{\hertz} e \SI{1000}{\hertz} e limitação de banda por um par de filtros de Butterworth de
terceira ordem \cite{iso2954_2012}.

O valor medido é comparado com quatro zonas de avaliação:

\begin{alineas}
  \item zona A: vibração típica de máquinas recém-comissionadas;
  \item zona B: vibração aceitável para operação de longo prazo sem restrição;
  \item zona C: vibração insatisfatória para operação contínua de longo prazo, em que a máquina pode operar por período limitado, até que haja
        oportunidade de ação corretiva;
  \item zona D: vibração suficientemente severa para causar dano à máquina.
\end{alineas}

Os limites entre as zonas dependem do porte da máquina e do tipo de fundação, conforme a Tabela~\ref{tab:fund-zonas}. A fundação é dita rígida
quando a primeira frequência natural do conjunto máquina--fundação está suficientemente acima da frequência de excitação principal, e flexível
no caso contrário. Os valores de zona são iguais nas edições de 2009 e de 2022.

\begin{table}[!htb]
\IBGEtab{%
\caption{Limites entre zonas de avaliação para máquinas industriais, em velocidade eficaz de banda larga}%
\label{tab:fund-zonas}%
}{%
\begin{tabular}{@{}llccc@{}}
\toprule
Grupo & Fundação & \multicolumn{3}{c}{Limite entre zonas (\si{\milli\meter\per\second})} \\
\cmidrule{3-5}
 & & A/B & B/C & C/D \\
\midrule
1 & rígida   & \num{2,3} & \num{4,5} & \num{7,1}  \\
1 & flexível & \num{3,5} & \num{7,1} & \num{11,0} \\
2 & rígida   & \num{1,4} & \num{2,8} & \num{4,5}  \\
2 & flexível & \num{2,3} & \num{4,5} & \num{7,1}  \\
\bottomrule
\end{tabular}%
}{%
\fonte{\citeonline{iso208163_2022}.}%
\nota{Grupo 1: potência acima de \SI{300}{\kilo\watt}, ou máquinas elétricas com altura de eixo de pelo menos \SI{315}{\milli\meter}. Grupo 2: potência de \SIrange{15}{300}{\kilo\watt}, ou máquinas elétricas com altura de eixo entre \SI{160}{\milli\meter} e \SI{315}{\milli\meter}.}%
}
\end{table}

\subsection{Integração de aceleração para velocidade}
\label{sec:fund-integracao}

Como a norma avalia a severidade em velocidade e o acelerômetro fornece aceleração, é preciso integrar o sinal. Há duas formas de fazer isso
\cite{brandt2014}.

No \emph{domínio do tempo}, a velocidade é aproximada pela soma acumulada das amostras de aceleração multiplicadas pelo intervalo de amostragem. O
método é simples, mas qualquer componente contínua presente na aceleração, como um desvio de zero do sensor, é integrada como uma rampa, e o erro
cresce linearmente com a duração do registro. Um filtro passa-alta aplicado antes da integração reduz esse efeito, mas introduz transitórios nas
bordas da janela.

No \emph{domínio da frequência}, usa-se a relação da Equação~\eqref{eq:fund-senoide}: cada componente de frequência $f$ da velocidade é a componente
correspondente da aceleração dividida por $\mathrm{j}2\pi f$. Calcula-se a DFT da aceleração, divide-se cada coeficiente por $\mathrm{j}2\pi f_k$,
anulam-se os coeficientes fora da banda de interesse e retorna-se ao tempo pela transformada inversa. Anular esses coeficientes elimina a componente
contínua e a divisão por zero em $f = 0$. \citeonline{brandt2014} compararam os dois métodos e concluíram que a integração no domínio da frequência é
a mais exata para registros longos. A busca não encontrou a aplicação desse método à velocidade eficaz de banda larga em um dispositivo de banda
limitada, ou seja, quanto um instrumento desse tipo subestima o indicador da norma.

Como exemplo, uma aceleração senoidal de amplitude \SI{3,0}{\meter\per\second\squared} a \SI{50}{\hertz} corresponde a uma velocidade de amplitude
$3{,}0/(2\pi \cdot 50) \approx \SI{9,55}{\milli\meter\per\second}$ e valor eficaz $\num{9,55}/\sqrt{2} \approx \SI{6,75}{\milli\meter\per\second}$. Pela
Tabela~\ref{tab:fund-zonas}, esse valor está na zona B para máquinas do grupo 1 em fundação flexível, na zona C para o grupo 1 rígido e o grupo 2
flexível, e na zona D para o grupo 2 rígido. Um erro de poucos milímetros por segundo muda, portanto, a classificação da máquina. A
seção~\ref{cap:resultados} usa esse valor analítico para comparar os métodos de integração.

%%%%%%%%%%%%%%%%%%%%%%%%%%%%%%%%%%%%%%%%%%%%%%%%%%%%%%%%%%%%%%%%%%%%%%%%%%%%%%%
\section{Detecção de anomalias e sua avaliação}
\label{sec:fund-anomalias}

Um detector de anomalias depende de três escolhas: o paradigma de aprendizado, o algoritmo que produz o escore e o limiar que transforma o escore em
decisão. O resultado informado também depende da métrica e da separação entre dados de treino e de teste.

\subsection{Paradigmas de aprendizado}

A condição da máquina pode ser avaliada de modos diferentes, conforme os dados disponíveis para treinar o modelo \cite{chandola2009}. Na abordagem
\emph{supervisionada}, o modelo aprende com exemplos rotulados de cada condição (normal e cada tipo de falha) e classifica novos dados em uma dessas
classes. Na abordagem \emph{semissupervisionada}, ou de classe única, o modelo aprende apenas com exemplos da condição normal e marca como anomalia o que
se afasta dela. Na abordagem \emph{não supervisionada}, não há rótulo algum, e o modelo supõe que as anomalias são raras e diferentes da maioria dos dados.

Em manutenção preditiva, a abordagem supervisionada esbarra em um problema prático: uma máquina recém-instalada não tem histórico de falhas, e reunir
exemplos rotulados de cada defeito exige que eles ocorram. Os dados de operação normal existem desde o primeiro dia. Por isso, um dispositivo instalado em
máquina sem histórico precisa treinar só com eles \cite{chandola2009}. Em contrapartida, o modelo detecta
\emph{afastamento} da normalidade e não identifica \emph{qual} defeito está presente.

\subsection{Isolation Forest}
\label{sec:fund-iforest}

O Isolation Forest \cite{liu2008} se baseia na ideia de que anomalias são poucas e diferentes e, por isso, mais fáceis de \emph{isolar} do resto dos dados do
que os pontos normais. O algoritmo constrói um conjunto de árvores de isolamento. Cada árvore usa uma subamostra aleatória dos dados de treino e a divide
recursivamente: em cada nó, escolhe-se ao acaso um atributo e um valor de corte entre o mínimo e o máximo desse atributo, até que cada ponto fique isolado
em uma folha ou que se atinja um limite de profundidade. Um ponto anômalo, por estar em região de baixa densidade, tende a ser isolado em poucas divisões,
perto da raiz. Um ponto normal exige mais divisões.

O comprimento do caminho $h(x)$ entre a raiz e a folha em que o ponto $x$ cai mede, portanto, quão fácil é isolá-lo. Para comparar árvores construídas com
subamostras de tamanho $\psi$, essa medida é normalizada pelo comprimento médio esperado do caminho em uma árvore binária de busca com $\psi$ pontos,

\begin{equation}
c(\psi) = 2H(\psi-1) - \frac{2(\psi-1)}{\psi}, \qquad H(i) \approx \ln i + \num{0,5772},
\label{eq:fund-cn}
\end{equation}

\noindent
e o escore de anomalia é

\begin{equation}
s(x, \psi) = 2^{-\,\mathbb{E}[h(x)] / c(\psi)},
\label{eq:fund-score}
\end{equation}

\noindent
em que $\mathbb{E}[h(x)]$ é a média do comprimento do caminho sobre todas as árvores \cite{liu2008}. O escore está entre 0 e 1. Valores próximos de 1 indicam
ponto fácil de isolar, candidato a anomalia. Valores bem abaixo de \num{0,5} indicam ponto normal. Se todos os pontos têm escore próximo de \num{0,5}, os dados
não contêm anomalias distintas. Os autores recomendam $t = 100$ árvores e subamostras de $\psi = 256$ pontos, para as quais $c(256) \approx \num{10,2}$.

Algumas características favorecem a execução embarcada: o treino usa subamostras pequenas, a inferência percorre cada árvore uma única vez com uma comparação
por nó, e o modelo treinado é inteiramente descrito por atributos e valores de corte \cite{liu2008}.

\subsection{Escolha do detector}
\label{sec:rev-detector}

Dois comparativos chegam a conclusões diferentes. O primeiro avalia \num{30} algoritmos de vários paradigmas
sobre \num{57} conjuntos de dados e conclui que \emph{nenhum} método não supervisionado é estatisticamente superior aos demais \cite{han2022}. O segundo compara
\num{33} algoritmos não supervisionados sobre \num{52} conjuntos tabulares reais e encontra vantagem para o Extended Isolation Forest, com a ressalva de que os
dados se separam em dois regimes \cite{bouman2024}.

Este trabalho usa o algoritmo original por causa do custo embarcado. Um nó da variante estendida guarda um vetor normal com tantas componentes quantos forem os
descritores, e a inferência passa de uma comparação para um produto interno por nó. \emandamento{Esse é um argumento de projeto. O consumo de memória e de tempo
ainda não foi medido no dispositivo.}

Para máquinas rotativas, o algoritmo original tem bons resultados publicados. Entre \num{11} detectores não supervisionados, ele atingiu F1 de \SI{97,19}{\percent}
em rolamento \emph{run-to-failure}, de \SI{99,71}{\percent} em caixa de engrenagens e de \SI{97,20}{\percent} em falhas mecânicas de bancada \cite{brito2022}. Os dois
primeiros conjuntos foram adquiridos a \SI{20}{\kilo\hertz}, banda que uma IMU de consumo não tem.

\subsection{Limiar de decisão e taxa de falso alarme}
\label{sec:fund-limiar}

O escore contínuo de anomalia precisa ser convertido em decisão por um \emph{limiar}: acima dele, o ponto é marcado como anômalo. A escolha do limiar troca um tipo de
erro pelo outro. Um limiar baixo detecta mais falhas e marca mais pontos normais como anomalia, e um limiar alto faz o contrário.

Na implementação de referência do algoritmo, a escolha é controlada pelo parâmetro de contaminação, que representa a fração esperada de anomalias nos dados de treino.
No valor padrão, o limiar corresponde à referência de \num{0,5} do escore, sem calibração para os dados. Com um valor numérico, o limiar é posicionado de modo que essa
fração dos pontos de treino fique acima dele. Nos dois casos, o parâmetro muda apenas um deslocamento fixo e não muda o escore de cada ponto \cite{liu2008,scikitlearn2026}.
Por isso, toda métrica de ordenação, inclusive a área sob a curva ROC (seção~\ref{sec:fund-roc}), não depende desse parâmetro.

Quando o treino contém apenas operação normal, o limiar pode ser calibrado diretamente pela taxa de falso alarme desejada: se o limiar é o quantil de ordem $1-\alpha$ dos
escores de dados normais não usados no ajuste das árvores, espera-se que uma fração $\alpha$ da operação normal fique acima dele. Outra alternativa modela apenas a cauda da
distribuição dos escores pela teoria de valores extremos e fixa o limiar com base em um risco alvo, isto é, na probabilidade de falso alarme admitida \cite{siffer2017}.

Os trabalhos que relatam bom desempenho não usam o valor padrão. Um fixa o escore normalizado em \num{0,75} sob a hipótese de treino limpo \cite{antonini2023}, e outro
calcula o limiar pela razão de contaminação conhecida \cite{brito2022}. As duas implementações embarcadas do mesmo detector que informam o falso alarme, ou permitem
calculá-lo, ficam entre \SI{22}{\percent} e \SI{38}{\percent} sobre operação normal \cite{chevtchenko2023,kolok2025}. Com essas taxas, um dispositivo emitiria de um a dois
alarmes falsos a cada cinco avaliações de uma máquina em operação normal.

\subsection{Matriz de confusão e taxas derivadas}

Com a anomalia como classe positiva, cada decisão do detector cai em uma de quatro categorias: verdadeiro positivo (VP), anomalia corretamente detectada; falso positivo (FP),
operação normal marcada como anomalia; verdadeiro negativo (VN); e falso negativo (FN), anomalia não detectada. Dessas contagens derivam a taxa de verdadeiros positivos (TPR),
a taxa de falsos positivos (FPR), a precisão e a medida $F_1$ \cite{fawcett2006}:

\begin{equation}
\mathrm{TPR} = \frac{\mathrm{VP}}{\mathrm{VP} + \mathrm{FN}},
\quad
\mathrm{FPR} = \frac{\mathrm{FP}}{\mathrm{FP} + \mathrm{VN}},
\quad
\text{precisão} = \frac{\mathrm{VP}}{\mathrm{VP} + \mathrm{FP}},
\quad
F_1 = \frac{2\cdot \text{precisão} \cdot \mathrm{TPR}}{\text{precisão} + \mathrm{TPR}}.
\label{eq:fund-taxas}
\end{equation}

A TPR é a \emph{taxa de detecção}, e a FPR é a \emph{taxa de falso alarme}. Em monitoramento contínuo, a taxa de falso alarme pesa na operação: um dispositivo que avalia a máquina a
cada poucos minutos e marca \SI{10}{\percent} da operação normal como anomalia gera dezenas de alarmes falsos por dia, e alarmes frequentes demais deixam de ser atendidos. A acurácia,
fração de decisões corretas, informa pouco quando as classes são desbalanceadas, porque um detector que nunca dispara alcança acurácia alta se as anomalias são raras.

\subsection{Curva ROC, área sob a curva e área parcial}
\label{sec:fund-roc}

Como o limiar é uma escolha, um detector tem uma TPR para cada limiar. A curva característica de operação do receptor (ROC) traça a TPR contra a FPR à medida que o limiar varre todos os
valores possíveis, e a área sob essa curva (AUC) resume o desempenho em um número entre \num{0,5}, o de uma decisão ao acaso, e 1, o de um detector perfeito \cite{fawcett2006}. A AUC é a
probabilidade de que uma anomalia escolhida ao acaso receba escore maior que um ponto normal escolhido ao acaso. Como depende apenas da ordem dos escores, ela não muda quando o limiar é
deslocado.

Um dispositivo de campo, porém, opera em um único ponto da curva, com falso alarme baixo, e dois detectores com a mesma AUC podem ter desempenhos muito diferentes nessa região. A
\emph{área parcial sob a curva} (pAUC) restringe a integração à faixa de FPR de interesse, entre 0 e um limite $\beta$, e a área nessa faixa é
$A_\beta = \int_{0}^{\beta} \mathrm{TPR}\; \mathrm{d}\,\mathrm{FPR}$. Este trabalho usa a forma padronizada de \citeonline{mcclish1989}, também calculada pela biblioteca usada no treino,

\begin{equation}
\mathrm{pAUC}(\beta) = \frac{1}{2}\left(1 + \frac{A_\beta - \beta^2/2}{\beta - \beta^2/2}\right),
\label{eq:fund-pauc}
\end{equation}

\noindent
que vale \num{0,5} para uma decisão ao acaso e 1 para um detector perfeito. O desafio de referência de detecção não supervisionada de anomalias em máquinas usa outra forma, $A_\beta/\beta$, com
$\beta = \num{0,1}$ \cite{koizumi2020}. Outro trabalho informa a área sob a curva de precisão e revocação (PR-AUC) além do F1 \cite{brito2022}.

\subsection{Validação cruzada e vazamento de informação}
\label{sec:fund-validacao}

O desempenho de um modelo precisa ser medido em dados não usados no treino. Na \emph{validação cruzada em $k$ partições}, os dados são divididos em $k$ partições, e o modelo é treinado $k$ vezes,
cada vez com uma partição reservada para teste. O desempenho é a média das $k$ avaliações, e sua dispersão indica quanto o resultado depende dos dados usados no teste.

A validação cruzada exige que os dados de treino e de teste sejam independentes. Quando não são, ocorre \emph{vazamento de informação}: o modelo tem acesso, no treino, a informação que não estaria
disponível na aplicação real, e o desempenho medido fica acima do verdadeiro \cite{kaufman2012}. Em sinais de vibração, a forma mais comum de vazamento vem da segmentação. Um ensaio de alguns
segundos é dividido em muitas janelas, e janelas do mesmo ensaio têm características que nada têm a ver com a falha, como nível de ruído, ganho do sensor e condição de fixação. Se as janelas são
sorteadas entre treino e teste, o modelo pode acertar porque reconhece o ensaio. \citeonline{varejao2025} chamam esse efeito de \emph{viés de similaridade}. A correção é a \emph{validação por
grupos}, aqui chamada de partição por ensaio: todas as janelas de um mesmo ensaio ficam na mesma partição, e o teste é sempre feito sobre ensaios que o modelo não viu.

Esse problema já foi publicado e não é contribuição deste trabalho. De \num{41} estudos revisados sobre uma mesma base pública de rolamentos, \num{40} usaram delineamento sujeito a esse vazamento
\cite{varejao2025}, e a diferença medida é da ordem de dezenas de pontos percentuais \cite{vieira2026}. Outras revisões sobre classificação de falha de rolamento chegam a conclusão parecida
\cite{abburi2023,smith2015}.

%%%%%%%%%%%%%%%%%%%%%%%%%%%%%%%%%%%%%%%%%%%%%%%%%%%%%%%%%%%%%%%%%%%%%%%%%%%%%%%
\section{Nós sensores e enlace LoRa}
\label{sec:fund-no}

O Kaelix é um nó sensor: um sistema embarcado que mede a vibração e envia o resultado por rádio.

\subsection{Sistemas embarcados para monitoramento}
\label{sec:fund-embarcados}

Um nó sensor de monitoramento é um sistema embarcado: um computador dedicado a uma única função, integrado ao equipamento que atende, com pouca memória e pouca capacidade de processamento. Seu
componente central é o microcontrolador, que reúne processador, memória e periféricos (conversores, interfaces seriais, temporizadores) em um único circuito integrado.

O Kaelix usa microcontroladores da família ESP32. O ESP32-S3, por exemplo, tem dois núcleos de até \SI{240}{\mega\hertz}, \SI{512}{\kilo\byte} de memória estática de acesso aleatório (SRAM) e
rádios Wi-Fi e Bluetooth integrados \cite{espressif2026}.

Executar a detecção no próprio dispositivo reduz o uso do rádio, porque transmitir uma janela de \num{512} amostras de três eixos ocuparia o enlace por muito mais tempo que transmitir o resultado. A
execução de modelos de aprendizado de máquina em microcontroladores, com memória de centenas de quilobytes e consumo de miliwatts, é o campo denominado TinyML \cite{antonini2023,barbariol2022}.

\subsection{Nós de baixo custo publicados}
\label{sec:rev-implantacao}

Os nós sem fio para manutenção preditiva dividem-se em dois grupos. Os produtos comerciais têm banda uma ordem de grandeza acima da de uma IMU de consumo
\cite{bentlynevada2025,erbessd2025,tractian2022}. Os nós acadêmicos de baixo custo usam a mesma arquitetura estudada aqui: microcontrolador, IMU MEMS com taxa na ordem de \SI{1}{\kilo\hertz} e
detecção não supervisionada \cite{antonini2023,bisio2025,kolok2025}.

O trabalho mais próximo em \emph{hardware} e em modo de falha relata acurácia de \SI{72,9 \pm 2,4}{\percent}, com validação cruzada de cinco partições sobre cerca de \num{900} amostras e \emph{um
único} modo de falha \cite{kolok2025}. O trabalho termina no valor eficaz de aceleração, sem integração para velocidade e sem critério da norma.

A comparação mais direta é embarcada e mede três detectores no mesmo nó e no mesmo conjunto \cite{chevtchenko2023}. Nela, um método baseado em vizinhança superou o algoritmo usado aqui em
especificidade e em latência, com execução aproximadamente \num{26} vezes mais rápida. \emandamento{O método de vizinhança precisa manter o conjunto de referência na memória, enquanto a floresta
guarda apenas limiares e índices, e esse foi o motivo para manter o detector atual. O consumo de memória no dispositivo ainda não foi medido.}

O mesmo trabalho mostra que um motor com cerca de um ano de uso tem assinatura normal diferente da de motores novos \cite{chevtchenko2023}. Um modelo único com limiar fixo sofre com essa mudança,
tanto entre máquinas quanto ao longo da vida útil.

Outros trabalhos mostram que o detector treina e executa dentro de um microcontrolador dessa classe \cite{antonini2023}, que há variante criada para reduzir seu custo em TinyML \cite{barbariol2022}
e que dispositivos da mesma família suportam processamentos mais pesados \cite{kumar2026,uselis2025}.

\subsection{Enlace LoRa}
\label{sec:fund-lora}

LoRa é uma técnica de modulação de rádio criada para enlaces de longo alcance e baixo consumo, com baixa taxa de dados \cite{augustin2016}. Ela usa espalhamento espectral por \emph{chirp}: cada
símbolo é representado por um sinal cuja frequência varia linearmente ao longo de toda a largura de banda do canal. O espalhamento permite demodular sinais abaixo do nível de ruído, o que aumenta o
alcance. LoRa designa a camada física, e LoRaWAN é um protocolo de rede definido sobre ela, com servidores de rede e gerência de dispositivos. Um enlace ponto a ponto, entre o dispositivo e um
receptor próprio, pode usar a modulação LoRa sem o protocolo LoRaWAN.

Três parâmetros definem o equilíbrio entre alcance, taxa e consumo \cite{semtech2015}:

\begin{alineas}
  \item fator de espalhamento (SF), de 6 a 12: cada símbolo carrega SF bits e é formado por $2^{\mathrm{SF}}$ \emph{chips}, e cada incremento de SF dobra a duração do símbolo e aumenta a
        sensibilidade do receptor;
  \item largura de banda (BW), tipicamente \SI{125}{\kilo\hertz}: banda mais larga encurta o símbolo e reduz a sensibilidade;
  \item taxa de codificação (CR), de 4/5 a 4/8: fração de redundância acrescentada para correção de erros.
\end{alineas}

A taxa de símbolos é $R_s = \mathrm{BW}/2^{\mathrm{SF}}$, e a duração do símbolo é $T_s = 1/R_s$. A duração total de um pacote, o \emph{tempo no ar}, é a soma da duração do preâmbulo e da carga útil
\cite{semtech2015}:

\begin{equation}
\begin{split}
T_{\mathrm{pacote}} = {} & (n_{\mathrm{pre}} + \num{4,25})\,T_s \\
& + \left[ 8 + \max\!\left( \left\lceil \frac{8\,\mathrm{PL} - 4\,\mathrm{SF} + 28 + 16\,\mathrm{CRC} - 20\,\mathrm{IH}}{4(\mathrm{SF} - 2\,\mathrm{DE})} \right\rceil (\mathrm{CR} + 4),\; 0 \right) \right] T_s,
\end{split}
\label{eq:fund-toa}
\end{equation}

\noindent
em que $n_{\mathrm{pre}}$ é o número de símbolos do preâmbulo, PL o número de bytes da carga útil, CRC vale 1 quando a verificação de redundância está habilitada, IH vale 1 quando o cabeçalho é
omitido, DE vale 1 quando a otimização para baixa taxa está ativa, e CR vale de 1 (para 4/5) a 4 (para 4/8).

O enlace do Kaelix usa SF9, BW~=~\SI{125}{\kilo\hertz}, CR~=~4/5, preâmbulo de 8 símbolos e carga útil de \num{21} bytes, com cabeçalho e CRC habilitados e sem otimização para baixa taxa. A duração
do símbolo é $2^9/\num{125000} = \SI{4,096}{\milli\second}$, o preâmbulo ocupa $(8 + \num{4,25}) \times \num{4,096} \approx \SI{50,2}{\milli\second}$, e a carga útil ocupa
$8 + \lceil 176/36 \rceil \times 5 = 33$ símbolos, ou \SI{135,2}{\milli\second}. O tempo no ar resultante é de aproximadamente \SI{185}{\milli\second}. Em SF12, a mesma carga útil ocuparia o rádio
por um tempo várias vezes maior, em troca de mais alcance.

%%%%%%%%%%%%%%%%%%%%%%%%%%%%%%%%%%%%%%%%%%%%%%%%%%%%%%%%%%%%%%%%%%%%%%%%%%%%%%%
\section{Síntese e lacunas}
\label{sec:rev-lacunas}

A literatura revisada tem mais resultados sobre aquisição, em que um sensor de consumo tem maiores limitações, do que sobre estimação, em que este trabalho concentra seus experimentos.

Os trabalhos revisados apoiam quatro pontos usados aqui. O limiar padrão não é calibrado, e há uma alternativa baseada na teoria de valores extremos. O \emph{benchmark} de referência informa a métrica
no ponto de operação. O detector já foi executado em microcontrolador. A norma exige a integração de aceleração para velocidade, e o método correto é conhecido.

Três escolhas deste trabalho vão contra o que a literatura recomenda. Nenhum estudo de adequação revisado testou sensor com limite superior próximo ao usado aqui, e a assinatura de rolamento depende de
uma banda que o dispositivo não adquire. O piso de ruído de uma IMU de consumo limita a medição, qualquer que seja o processamento. Por fim, a maior comparação de métodos não supervisionados recomenda
outra variante do detector, e um método de vizinhança foi mais específico e mais rápido no mesmo tipo de \emph{hardware} embarcado. \emandamento{Este trabalho ainda não mediu o consumo de memória e de
tempo no dispositivo, que é o argumento para manter o detector atual.}

A revisão encontrou sete lacunas. Duas lacunas supostas no início da revisão não se confirmaram: dois trabalhos informam o falso alarme do
detector em nó embarcado, ou permitem calculá-lo \cite{chevtchenko2023,kolok2025}, e, com isso, a fração da operação normal marcada pelo limiar padrão já está quantificada.

\textbf{Lacuna 1: efeito isolado do parâmetro de contaminação.} Não se encontrou trabalho que variasse o parâmetro sobre a \emph{mesma} base e o \emph{mesmo} modelo e informasse a fração da operação normal
        marcada.

\textbf{Lacuna 2: falso alarme de detector treinado no próprio dispositivo.} O único trabalho revisado que treina no próprio dispositivo deixa a avaliação de acurácia fora do escopo \cite{antonini2023}.

\textbf{Lacuna 3: contraste diagnóstico em função da banda.} A literatura discute o efeito da banda de forma qualitativa, e o único estudo revisado que varre a taxa de amostragem mede a acurácia de um
        classificador supervisionado, e não o contraste do indicador \cite{elbouharrouti2024}.

\textbf{Lacuna 4: rebatimento sobre descritores estatísticos.} O efeito não foi encontrado em literatura revisada por pares e aparece apenas em notas de aplicação de fabricante.

\textbf{Lacuna 5: efeito da limitação de banda sobre o indicador da norma.} O erro dos métodos de integração está quantificado \cite{brandt2014}, mas não sua aplicação a esse indicador em dispositivo de banda
        limitada.

\textbf{Lacuna 6: detecção limitada pelo piso de ruído.} Não se encontrou análise de defeito mínimo detectável em função da densidade de ruído para IMUs de consumo.

\textbf{Lacuna 7: separação entre efeito de banda e efeito de montagem.} Não se encontrou \emph{benchmark} aberto com sinais gravados \emph{simultaneamente} por MEMS de banda limitada e por piezoelétrico no
        mesmo ponto.

A seção~\ref{cap:resultados} trata das lacunas 1 e 4 em sinal sintetizado e oferece à lacuna 3 apenas dois pontos, insuficientes para formar uma função da banda. As demais continuam abertas, e parte delas
aparece como trabalho futuro na seção~\ref{sec:conc-futuro}.
